\scanpage{113}
% archive foliation: 112

alors qu'il est clair que $z_0=\deleted{\unsure{x_0}}\inserted{\unsure{x_0}},\ z_0=y_0$
sont solutions solutions de l'équation (48).
Si donc $z_0$ est une solution $\ne x_0,y_0$, on aura
\[
\begin{aligned}
x_0+y_0+z_0&=-(a_1-\lambda^2)=\lambda^2-a_1\\
&=\left(\frac{x'-y'}{x_0-y_0}\right)^2\unsure{-}a_1
\end{aligned}
\]
d'où la solution (sous réserve d'existence d'une solution $z_0\ne x_0,y_0$,
ou encore \deleted{\unsure{d'un $z$}} \inserted{réserve} que
\[
(50)\quad\left\{\begin{aligned}
&x,y\in E\setminus\{0\},\quad x\ne y,\quad -(x+y)\ne x,y\quad\text{i.e.}\\
&\hspace{3cm}2x+y\ne0,\quad2y+x\ne0\\
&\hspace{3cm}\text{i.e. }y\ne-2x,\ x\ne-2y
\end{aligned}\right.
\]
), la solution
\[
(51)\quad\left\{\begin{aligned}
z_0&=-x_0-y_0+\left(\frac{x'-y'}{x_0-y_0}\right)^2-a_1\\
z'&=\lambda z_0+\mu=\frac{x'-y'}{x_0-y_0}
\left[-x_0-y_0+\left(\frac{x'-y'}{x_0-y_0}\right)^2-a_1\right]
\inserted{{}+\frac{y'x_0-x'y_0}{x_0-y_0}}
\end{aligned}\right.
\]

Par principe de prolongement, cette expression pour $z_0=\wp(z)$ et $z'=\wp'(z)$,
pour $z=-(x+y)$, est valable encore dès que $x,y\in E(k)\setminus\{0\}$
sont tels que $x_0\ne y_0$ i.e. $\wp(x)\ne\wp(y)$, i.e. dès que
\deleted{$\unsure{x=}$} $y\notin\pm x$, de sorte que les formules (51) ont un sens.
